\subsection{MODULES; VECTOR SPACES; LINEAR COMBINATIONS}

DEFINITION 1. Given a ring A, a left module over A (or left A-module) is a set E with an algebraic structure defined by giving:

\begin{enumerate}
\def\labelenumi{(\arabic{enumi})}
\item
  a commutative group law on E (written additively in what follows);
\item
  a law of action \((\alpha, x) \mapsto \alpha \tau x\) , whose domain of operators is the ring A and which satisfies the following axioms:
\end{enumerate}

\((\mathbf{M}_{\mathbf{I}})\alpha \tau (x + y) = (\alpha \tau x) + (\alpha \tau y)\) for all \(\alpha \in \mathbf{A}, x \in \mathbf{E}, y \in \mathbf{E}\) ;

\((M_{\mathrm{II}}) (\alpha + \beta) \tau x = (\alpha \tau x) + (\beta \tau x) \text{ for all } \alpha \in A, \beta \in A, x \in E;\)

\((M_{\mathrm{III}})\alpha \tau (\beta \tau x) = (\alpha \beta)\tau x\) for all \(\alpha \in A,\beta \in A,x\in E;\)

\((M_{\mathrm{IV}}) \ 1 \tau x = x \ for \ all \ x \in E.\)

Axiom (M \(_{I}\) ) means that the external law of a left A-module E is distributive with respect to addition on E; a module is thus a commutative group with operators.

If in Definition 1, axiom \((\mathbf{M}_{\mathrm{III}})\) is replaced by

\[
\left(\mathrm{M} _ {\mathrm{III}} ^ {\prime}\right) \alpha \tau (\beta \tau x) = (\beta \alpha) \tau x for all \alpha \in \mathrm{A}, \beta \in \mathrm{A}, x \in \mathrm{E},
\]

E with the algebraic structure thus defined is a right module over A or a right A-module.

When speaking of A-modules (left or right), the elements of the ring A are often called scalars.

Usually the external law of composition of a left A-module (resp. right A-module) is written multiplicatively, with the operator written on the left (resp. right); condition \((\mathbf{M}_{\mathrm{III}})\) is then written \(\alpha(\beta x) = (\alpha\beta)x\) , condition \((\mathbf{M}_{\mathrm{III}}^{\prime})\) is written \((x\beta)\alpha = x(\beta\alpha)\) .

If \(A^{0}\) denotes the opposite ring to A (I, §8, no. 3), every right module E over the ring A is a left module over the ring \(A^{0}\) . It follows that an exposition can be given of the properties of modules whilst systematically confining attention either to left modules or to right modules; in §§ 1 and 2 we shall in general give this exposition for left modules and when we speak of a module (without specifying which type) we shall mean a left module whose external law will be written multiplicatively. When the ring A is commutative the notions of right module and left module with respect to A are identical.

For all \(\alpha \in \mathbf{A}\) , the mapping \(x \mapsto \alpha x\) of an A-module \(\mathbf{E}\) into itself is called the homothety of ratio \(\alpha\) of \(\mathbf{E}\) (I, §4, no. 2); by \((\mathbf{M}_{\mathbf{I}})\) a homothety is an endomorphism of the commutative group structure (without operators) on \(\mathbf{E}\) , but not in general of the module structure on \(\mathbf{E}\) (I, §4, no. 2; cf.~II, §1, no. 2 and no. 13). Hence \(\alpha 0 = 0\) and \(\alpha(-x) = -(\alpha x)\) ; if \(\alpha\) is an invertible element of \(\mathbf{A}\) , the homothety \(x \mapsto \alpha x\) is an automorphism of the commutative group structure (without operators) on \(\mathbf{E}\) , for the relation \(y = \alpha x\) implies by virtue of \((\mathbf{M}_{\mathbf{IV}})\) , \(x = \alpha^{-1}(\alpha x) = \alpha^{-1}y\) .

Similarly, by virtue of \((\mathbf{M}_{\mathrm{II}})\) , for all \(x \in E\) , the mapping \(\alpha \mapsto \alpha x\) is a homomorphism of the additive group A into the commutative group (without operators) E; hence 0x = 0 and \((- \alpha)x = -(\alpha x)\) ; moreover, by \((\mathbf{M}_{\mathrm{IV}})\) , for every integer \(n \in Z\) , \(n.x = (n.1)x\) .

When the ring A consists only of the element 0, every A-module E consists only of the element 0, for then \(1 = 0\) in A, whence, for all \(x \in \mathbf{E}\) , \(x = 1.x = 0.x = 0\) .

Examples. (1) Let \(\phi\) be a homomorphism of a ring A into a ring B; the mapping \((a, x) \mapsto \phi(a)x\) (resp. \((a, x) \mapsto x\phi(a)\) ) of A × B into B defines on B a left (resp. right) A-module structure. In particular if \(\phi\) is taken to be the identity mapping on A, a canonical left (resp. right) A-module structure is obtained on A; to avoid confusion, the set A with this structure is denoted by A \(_s\) (resp. A \(_d\) ).

\begin{enumerate}
\def\labelenumi{(\arabic{enumi})}
\setcounter{enumi}{1}
\item
  On a commutative group G (written additively) the group with operators structure defined by the external law \((n, x) \mapsto n.x\) (I, §3, no. 1) is a module structure over the ring Z of rational integers.
\item
  Let E be a commutative group written additively, \(\mathcal{E}\) the endomorphism ring of E (I, § 8, no. 3: recall that the product fg of two endomorphisms is by definition the composite endomorphism \(f \circ g\) ). The external law \((f, x) \mapsto f(x)\) between operators \(f \in \mathcal{E}\) and elements \(x \in E\) defines on E a canonical left \(\mathcal{E}\)-module structure.
\end{enumerate}

Consider now a ring A and suppose there is given on E a left (resp. right) A-module structure; for all \(\alpha \in \mathbf{A}\) , the homothety \(h_{\alpha}: x \mapsto \alpha x\) (resp. \(x \mapsto x\alpha\) ) belongs to \(\mathcal{E}\) ; the mapping \(\phi: \alpha \mapsto h_{\alpha}\) is a homomorphism of the ring A (resp. the opposite ring \(\mathbf{A}^0\) ) into the ring \(\mathcal{E}\) and by definition \(\alpha x = (\phi(\alpha))(x)\) (resp. \(x\alpha = (\phi(\alpha))(x)\) ). Conversely, giving a ring homomorphism \(\phi: \mathbf{A} \to \mathcal{E}\) (resp. \(\phi: \mathbf{A}^0 \to \mathcal{E}\) ) defines a left (resp. right) A-module structure on E by the above formulae. In other words, being given a left (resp. right) A-module structure on an additive group E with additive law the given group law is equivalent to being given a ring homomorphism \(\mathbf{A} \to \mathcal{E}\) (resp. \(\mathbf{A}^0 \to \mathcal{E}\) ).

DEFINITION 2. A left (resp. right) vector space over a field K is a left (resp. right) K-module.

The elements of a vector space are sometimes called vectors.

Examples. (4) A field is both a left and a right vector space with respect to any of its subfields.

*(5) The real number space of three dimensions \(\mathbf{R}^3\) is a vector space with respect to the field of real numbers \(\mathbf{R}\) , the product \(tx\) of a real number \(t\) and a point \(x\) with coordinates \(x_1, x_2, x_3\) being the point with coordinates \(tx_1, tx_2, tx_3\) . Similarly, the set of real-valued functions defined on an arbitrary set \(\mathbf{F}\) is a vector space with respect to \(\mathbf{R}\) , the product \(tf\) of a real number \(t\) and a function \(f\) being the real-valued function \(x \mapsto tf(x)\).*

For two families \((x_{i})_{i\in I}, (y_{i})_{i\in I}\) of elements of an A-module E of finite support (I, §2, no. 1), the following equations hold:

\begin{enumerate}
\def\labelenumi{(\arabic{enumi})}
\tightlist
\item
\end{enumerate}

\[
\sum_ {i \in I} (x _ {i} + y _ {i}) = \sum_ {i \in I} x _ {i} + \sum_ {i \in I} y _ {i}\tag{2}
\]

\[
\alpha . \sum_ {i \in I} x _ {i} = \sum_ {i \in I} (\alpha x _ {i}) \quad \text { for   all } \alpha \in A;
\]

the equations are immediately reduced to the analogous equations for finite sums by considering a finite subset H of I containing the supports of \((x_{i})\) and \((y_{i})\) .

DEFINITION 3. An element \(x\) of an A-module \(\mathbf{E}\) is said to be a linear combination with coefficients in \(\mathbf{A}\) of a family \((a_{i})_{i\in I}\) of elements of \(\mathbf{E}\) if there exists a family \((\lambda_{i})_{i\in I}\) of elements of \(\mathbf{A}\) , of finite support, such that \(x = \sum_{i\in I}\lambda_{i}a_{i}\) .

In general there are several distinct families \((\lambda_{i})\) satisfying this condition (cf.~no. 11).

Note that 0 is the only linear combination of the empty family of elements of E (by the convention of I, § 2, no. 1).

